\documentclass[12pt]{article}

\usepackage[a4paper, margin=2.5cm]{geometry}
\usepackage{amsmath}
\usepackage{amssymb}
\usepackage{amsthm}
\usepackage[colorlinks=true, linkcolor=blue, citecolor=blue, urlcolor=blue]{hyperref}
\usepackage{booktabs}
\usepackage{natbib}

\newtheorem{theorem}{Theorem}
\newtheorem{proposition}[theorem]{Proposition}
\newtheorem{corollary}[theorem]{Corollary}
\newtheorem{lemma}[theorem]{Lemma}
\theoremstyle{definition}
\newtheorem{definition}[theorem]{Definition}
\newtheorem{axiom}[theorem]{Axiom}
\newtheorem{hypothesis}[theorem]{Hypothesis}
\theoremstyle{remark}
\newtheorem{remark}[theorem]{Remark}

\newcommand{\R}{\mathbb{R}}
\newcommand{\C}{\mathbb{C}}
\newcommand{\Q}{\mathbb{Q}}
\newcommand{\Z}{\mathbb{Z}}
\newcommand{\Zphi}{\mathbb{Z}[\varphi]}
\newcommand{\twoI}{2\mathbf{I}}

\title{A Structural Correspondence for the Fine-Structure Constant\\[6pt]
\large $\alpha^{-1} \approx 137 + \pi/87$ under $F1$, {C1} (with C2, C3 as ambient narrative), {A}, {B}, {Kostant}, and {P2}}
\author{%
  Lee Smart\\[4pt]
  \textit{Institute of Vibrational Field Dynamics}\\[2pt]
  \texttt{contact@vibrationalfielddynamics.org}\\[2pt]
  \texttt{@vfd\_org}%
}
\date{April 2026}

\begin{document}

\maketitle

\begin{abstract}
We record the fine-structure constant correspondence $\alpha^{-1} \approx 137 + \pi/87$ conditional on the permeability axiom $F1$ ($r = 1 + 1/r$, forcing $r = \varphi$) together with four downstream ingredients already present in the Vibrational Field Dynamics programme: (i) the Elser--Sloane realisation of $E_8$ as the diagonally-embedded sublattice $\{(x, \sigma(x)) : x \in \mathcal{I}\} \subset \mathcal{I} \oplus \sigma(\mathcal{I})$, where $\mathcal{I}$ is the icosian ring and $\sigma$ is the Galois involution (classical); (ii) the 12-dimensional ambient $L_{12} = E_8 \oplus \Zphi^2$ with C1 (phason-complement uniqueness) as an explicit hypothesis; C2 (upward termination) and C3 ($\Zphi$ self-reference) are part of the ambient 12D narrative but are \emph{not} used in the final $137 + \pi/87$ derivation proper; (iii) two structural assumptions A and B on the $\Zphi^2$-phason alignment with the four $E_8$ Coxeter planes and the role of the upper Coxeter exponents $\{17, 19, 23, 29\}$; and (iv) the $2I$-isotypic substrate on the 600-cell vertex graph imported from Paper~P2. Under these inputs, the integer $137$ decomposes as $87 + 50$ where $87$ is the quantity counted as the phason-sector slot total under Hypothesis~B and the gauge convention of Hypothesis~Kostant (subject to one global $U(1) \times \Z/2$ gauge reduction, which this paper treats as a convention, not a derivation), and $50 = 9 + 12 + 14 + 15$ is the sum of the four nontrivial $\Q$-rational $2I$-isotypic Laplacian eigenvalues of the 600-cell. The irrational $\pi/87$ remainder is recorded as a structural identification between Coxeter phase, Galois halving, and the slot count; this paper does \emph{not} claim a theorem-level derivation of $\pi/87$. The agreement with CODATA~2018 is $0.81$ ppm: $\alpha^{-1}_{\text{CODATA}} = 137.035\,999\,084(21)$ vs.\ $\alpha^{-1}_{\text{identity}} = 137.036\,110\,260\ldots$

\vspace{4pt}
\noindent\textbf{Epistemic status.} This paper is a \emph{hypothesis-audited structural identification}, not a theorem chain from $F1$. Five distinct open items are named explicitly (B-formalisation, A-canonicality, C1 minimality, Kostant gauge-fixing, and the irrational-term derivation); each is a polish gap between the current derivation and a first-principles proof.
\end{abstract}

\section{Introduction}\label{sec:intro}

This paper records the structural correspondence
\begin{equation}\label{eq:main}
  \boxed{\; \alpha^{-1} \;\approx\; 137 + \frac{\pi}{87} \;}
\end{equation}
as a conditional identification under the permeability axiom $F1$~\citep{CascadeFoundations}:
\begin{equation}\label{eq:F1}
  r = 1 + \frac{1}{r},
\end{equation}
which has positive solution $r = \varphi = (1 + \sqrt 5)/2$. Identity~\eqref{eq:main} was first reported as a numerical correspondence in~\citep{SmartV, SmartXXII}; the present paper collects the hypothesis-audited structural chain
\[
  F1 \;\longrightarrow\; \Zphi \;\longrightarrow\; \mathcal{I}
  \;\longrightarrow\; E_8 = \{(x,\sigma(x)):x\in\mathcal{I}\}
  \;\longrightarrow\; L_{12} = E_8 \oplus \Zphi^2

  \;\longrightarrow\; H_4 \text{ (cut-and-project)}
  \;\longrightarrow\; 87 + 50 + \pi/87.
\]

\paragraph{What is derived, what is imported, what remains open.}
The chain splits into four layers. For each, we summarise the epistemic status, cite the establishing source, and state what remains to be done before each layer has textbook rigour.

\begin{itemize}
  \item[(L-F1)] \emph{$F1$ forces $\varphi$ and $\Zphi$}, and $\Zphi$ is the minimal $\Z$-ring containing $\varphi$. Classical.
  \item[(L-E8)] \emph{$E_8 = \{(x,\sigma(x)) : x \in \mathcal{I}\} \subset \mathcal{I} \oplus \sigma(\mathcal{I})$ via the icosian ring}, after Elser--Sloane. Classical~\citep{ElserSloane1987, MoodyPatera1993}.
  \item[(L-12D)] \emph{$L_{12} = E_8 \oplus \Zphi^2$}, with its 12-dimensional closure and self-referential closure of the $\Zphi$ base. Conjectures C1--C3 in~\citep{Cascade12D}. C2 is argued in the admissible-ambient framework of \texttt{cascade-12d-closure.tex}~\citep{Cascade12D}; the earlier Phase~B markdown note~\citep{CascadePhaseB} gives an informal dimensional-counting precursor. C1 closes only under the $H_{\min}$ minimality principle together with the admissibility hypotheses catalogued in~\citep{Cascade12D}; C3 then follows conditionally on C1 plus classical $\Zphi$ structure.
  \item[(L-AB)] \emph{Canonical $\Zphi^2$-alignment with the four $E_8$ Coxeter planes; phason role of the upper Coxeter exponents $\{17, 19, 23, 29\}$.} Assumptions A and B in~\citep{CascadeABclosure}. Both are formulated in an untranslated markdown note (written in the $\Zphi^4$ convention); this paper treats them as explicit hypotheses, not as downstream-ready proofs. A closes modulo the Coxeter-plane-canonicality claim.
  \item[(L-P2)] \emph{600-cell vertex graph: isotypic $94{+}26$ decomposition and four nontrivial $\Q$-rational Laplacian eigenvalues $\{9, 12, 14, 15\}$ summing to $50$.} Recorded in P2~\citep{CascadeAlgSubstrate}, \S 6, via exact shell/class computations in $\Q(\sqrt 5)$ together with an imported $2I$ character table from Du~Val~\citep{DuVal1964}.
\end{itemize}

Identity~\eqref{eq:main} is obtained conditional on the full hypothesis load $\{F1, \text{C1 minimality}, \text{A canonicality}, B, \text{Kostant-gauge fixing}, \text{P2 character-table import}, \text{irrational-term structural identification}\}$. Under those conditions, the \emph{integer} part $137 = 87 + 50$ is traced structurally as below. The \emph{irrational} part $\pi/87$ is recorded as a structural identification with four unresolved heuristic steps (Remark~\ref{rem:irrational} and Section~\ref{sec:status}).

\paragraph{Structure of the paper.} Section~\ref{sec:F1} records $F1$ and its immediate consequences. Section~\ref{sec:L12} restates C1--C3 and notes the closures achieved in Phase~B. Section~\ref{sec:AB} states assumptions A and B and their role. Section~\ref{sec:P2} imports the 600-cell spectral substrate from P2. Section~\ref{sec:theorem} assembles the $\alpha$-structural correspondence. Section~\ref{sec:numerical} compares to CODATA~2018. Section~\ref{sec:status} audits the hypothesis load.

\section{The permeability axiom $F1$ and its immediate consequences}\label{sec:F1}

\begin{axiom}[$F1$, permeability]\label{axiom:F1}
The dimensionless permeability $r \in \R_{>0}$ of the discrete cascade substrate satisfies
\[
  r = 1 + \frac{1}{r}.
\]
\end{axiom}

The positive solution is $r = \varphi = (1 + \sqrt 5)/2$; hence all downstream ratios and ring-theoretic objects live in the unique minimal $\Z$-ring extension containing $\varphi$, namely $\Zphi$ (the ring of integers of $\Q(\sqrt 5)$, a PID with class number $1$).

\begin{lemma}[Elser--Sloane realisation of $E_8$]\label{lem:ES}
Let $\mathcal{I} \subset \mathbb{H}$ be the icosian ring: the quaternionic order of rank $4$ over $\Zphi$ (equivalently, rank $8$ over $\Z$). Let $\sigma: \Zphi \to \Zphi$ be the Galois involution $\sigma(\varphi) = 1 - \varphi = -1/\varphi$, extended coefficientwise. Then the $\Z$-module $\{(x, \sigma(x)) : x \in \mathcal{I}\} \subset \R^4 \oplus \R^4$ is isometric to the $E_8$ root lattice, and the cut-and-project $\pi_H: \mathcal{I} \to \R^4$ with golden-ratio acceptance window recovers the $H_4$ quasicrystal.
\end{lemma}

\begin{proof}
Classical; see~\citep[\S 3]{ElserSloane1987} and~\citep{MoodyPatera1993}. The icosian-ring construction of $E_8$ is collected in~\citep[\S 3--4]{CascadeAlgSubstrate} as P2~\S 3--4, with explicit coordinates for the projection $\pi_H: \Lambda(E_8) \to H_4$.
\end{proof}

\section{The 12-dimensional closure $L_{12}$ (conjectures C1--C3, Phase~B)}\label{sec:L12}

The upward cascade posits a 12-dimensional ambient lattice $L_{12} \subset \R^{12}$ containing $E_8 \subset \R^8$ together with a rank-$4$ phason complement $M$, i.e.\
\begin{equation}\label{eq:L12decomp}
  L_{12} = E_8 \oplus M
\end{equation}
as a $\Z$-module direct sum. Three conjectures gate the downstream $\alpha$-chain.

\begin{hypothesis}[C1, phason-complement uniqueness]
The rank-$4$ $\Z$-module $M$ in~\eqref{eq:L12decomp} is isomorphic to $\Zphi^2$ as a $\Zphi$-module.
\end{hypothesis}

\begin{hypothesis}[C2, upward termination at $L_{12}$]
No proper $\Z$-lattice strictly containing $L_{12}$ carries a further rank-$\geq 1$ cut-and-project extension of the type used in Lemma~\ref{lem:ES}.
\end{hypothesis}

\begin{hypothesis}[C3, $\Zphi$ self-reference]
The ring $\Zphi$ simultaneously (i) generates the substrate ring (Axiom~\ref{axiom:F1}), (ii) equals the coefficient ring of the icosian quaternionic order $\mathcal{I}$, and (iii) equals the scalar ring of the phason complement $M$.
\end{hypothesis}

\begin{remark}[Phase~B status of C1--C3]
Phase~B~\citep{CascadePhaseB} argues C2 by a bespoke dimensional-counting argument (not a classical theorem of algebraic number theory); argues C3 conditional on C1 plus classical $\Zphi$ structure; and argues C1 modulo a minimality principle $H_{\min}$ requiring $M$ to carry no ring structure strictly larger than $\Zphi$. The C1 argument additionally assumes $M$ is torsion-free and $\Zphi$-linear, neither of which is proved here; both are carried as implicit hypotheses of the C1 closure.
\end{remark}

\begin{remark}[Notation reconciliation: $\Zphi^2$ vs $\Zphi^4$]
\label{rem:notation}
The $\Zphi^2$ and $\Zphi^4$ readings are \emph{asserted in local sources} \citep[\S 8]{Cascade12D} to be alternative readings of the same underlying algebraic object, via a ring-rank relabelling ($\Zphi^2$ uses $\Zphi$-rank with $\mathrm{rank}_\Z M = 4$, $\mathrm{rank}_{\Zphi} M = 2$; $\Zphi^4$ uses Z-rank-style shorthand with $\mathrm{rank}_\Z M = 4$). We do not claim to \emph{exhibit} a reusable equivalence here: no explicit rewrite transporting the A, B, and $\pi/87$ arguments of the markdown notes~\citep{CascadePhaseB, CascadeABclosure} from their $\Zphi^4$ presentation into the $\Zphi^2$ convention has been carried out in the repository. The present paper uses the $\Zphi^2$ convention throughout; downstream citation of the markdown-note content requires such a translation, and its execution is part of the A/B polish agenda of \S\ref{sec:status}.
\end{remark}

\section{Assumptions A and B: Coxeter plane alignment and phason exponents}\label{sec:AB}

\begin{hypothesis}[A, $\Zphi^2$--Coxeter-plane alignment]\label{hyp:A}
There exist four mutually orthogonal $2$-planes $P_1, P_2, P_3, P_4 \subset \R^8$, each invariant under the $E_8$ Coxeter element $w$, such that the $\Zphi^2$ phason complement of C1 aligns canonically with the direct sum $P_1 \oplus P_2 \oplus P_3 \oplus P_4$ under the cut-and-project $L_{12} \to E_8$.
\end{hypothesis}

\begin{hypothesis}[B, upper Coxeter exponents as phason windings]\label{hyp:B}
The eigenvalues of $w$ on the four Coxeter planes $P_1, \ldots, P_4$ are $\exp(2\pi i m/30)$ for $m \in \{17, 19, 23, 29\}$ (the four \emph{upper} Coxeter exponents of $E_8$). Each upper exponent parameterises the phason winding along the corresponding $P_k$ per Coxeter cycle.
\end{hypothesis}

\begin{remark}
The $E_8$ Coxeter element $w$ has order $30$ and eigenvalues $\exp(2\pi i m/30)$ for $m \in \{1, 7, 11, 13, 17, 19, 23, 29\}$ (Bourbaki--Humphreys). Assumption B is the content of the ``lower/upper exponent'' dichotomy registered in~\citep{SmartXXII}: the four lower exponents $\{1, 7, 11, 13\}$ describe the physical (vertical) half, and the four upper exponents $\{17, 19, 23, 29\}$ describe the phason (internal) half under the Galois involution $\sigma$.
\end{remark}

\begin{hypothesis}[Kostant gauge reduction]\label{hyp:kostant}
A single global $U(1) \times \Z/2$ gauge orbit --- one dimensionless choice of origin shared across the four phason directions --- removes exactly one phason slot from the Coxeter-exponent-sum total.
\end{hypothesis}

\begin{proposition}[Phason slot count, as a counting convention]\label{prop:87}
Under Hypothesis~\ref{hyp:B} (upper Coxeter exponents identified with phason-direction windings) and Hypothesis~\ref{hyp:kostant} (Kostant gauge reduction), the paper assigns the count $87$ to the phason sector by summing the upper exponents and imposing the one-slot gauge reduction:
\[
  (17 + 19 + 23 + 29) \;-\; 1 \;=\; 88 \;-\; 1 \;=\; 87.
\]
Calling this quantity a ``phason-slot count'' presupposes the additional identification (part of Hypothesis~\ref{hyp:B}) that the summed upper-exponent winding equals the phason-slot count in the first place; that identification is adopted here as a counting convention, not derived. Hypothesis~\ref{hyp:A} (A canonicality) is used only to \emph{interpret} the four phason directions as aligned with four orthogonal Coxeter planes; it does not enter the arithmetic of this proposition.
\end{proposition}

\begin{proof}
The arithmetic $17 + 19 + 23 + 29 = 88$ is elementary given B. No theorem in this paper or in any cited local source derives the reduction $88 \to 87$; Hypothesis~\ref{hyp:kostant} is a convention only. The Kostant~1959 principal-orbit counting~\citep{Kostant1959} is consistent with a one-slot reduction but does not single-handedly select the shared global origin used here.
\end{proof}

\section{The P2 isotypic substrate}\label{sec:P2}

\begin{proposition}[Rational-irrep eigenvalue sum, from P2]\label{prop:50}
Let $L$ be the Laplacian of the $600$-cell vertex graph, and let $\{0, 9, 12, 14, 15\}$ be its five $\Q$-rational eigenvalues (the eigenvalues attached to the trivial, $\mathbf{4}$-, $\mathbf{5}$-, $\mathbf{6}$-, $\mathbf{4'}$-isotypic components of $\C[\twoI]$ in the P2 decomposition~\citep{CascadeAlgSubstrate}, \S 6, Fact on the Laplacian spectrum). The sum of the four nontrivial $\Q$-rational eigenvalues is
\[
  9 + 12 + 14 + 15 \;=\; 50.
\]
\end{proposition}

\begin{proof}
Immediate from P2~\S 6, Table on the complete Laplacian spectrum of the $600$-cell vertex graph. The isotypic decomposition $\C^{V_{600}} = W_{\mathrm{int}} \oplus W_{\mathrm{irr}}$ with $\dim W_{\mathrm{int}} = 94$ and $\dim W_{\mathrm{irr}} = 26$ was also stated as Theorem~1 of~\citet{SmartXXII}; both papers use the same $2I$ character-table import from~\citep{DuVal1964}.
\end{proof}

\section{The $\alpha$-structural correspondence}\label{sec:theorem}

\begin{proposition}[Integer part of the structural identity, conditional]\label{thm:alpha}
Under Hypotheses~\ref{hyp:B} and~\ref{hyp:kostant} (used via Proposition~\ref{prop:87}) and the P2 character-table import (used via Proposition~\ref{prop:50}),
\begin{equation}
  137 \;=\; 87 + 50
\end{equation}
where $87$ is the quantity counted as the phason-sector slot total under Hypothesis~\ref{hyp:B} and the gauge convention of Hypothesis~\ref{hyp:kostant}, and $50$ is the sum of the four nontrivial $\Q$-rational Laplacian eigenvalues of the 600-cell. The ambient hypotheses $F1$, C1--C3, and A enter this paper at earlier layers (they select the ring $\Zphi$, the phason-complement rank, and the Coxeter-plane interpretation respectively) but do not enter the arithmetic of this proposition.
\end{proposition}

\begin{proof}
By Propositions~\ref{prop:87} and~\ref{prop:50},
\[
  (17 + 19 + 23 + 29) - 1 \;+\; (9 + 12 + 14 + 15) \;=\; 87 \;+\; 50 \;=\; 137.
\]
\end{proof}

\begin{remark}[Irrational part $\pi / 87$: unproved heuristic identification]\label{rem:irrational}
The irrational $\pi/87$ remainder of Identity~\eqref{eq:main} is \emph{not} derived in this paper, and no cited source in the VFD repository contains a closed proof of it. The heuristic identification in~\citep{CascadeIdentityC, CascadeABclosure} proceeds by (i) assigning $\pi$ of Coxeter phase to each upper exponent, summing to $88\pi$; (ii) a two-to-one Galois operation on the $E_8 \to H_4$ projection (which, at the lattice-set level, projects $240 \to 120$ minimal vectors per P2 \S 4~\citep{CascadeAlgSubstrate}, but whose action on Coxeter phase is not proved in P2); (iii) a $\Z/2$ centre absorption of a $\pi$ residue; followed by (iv) division by the physical phason slot count $87$ to give $\pi/87$ as a per-slot residual. None of steps (i), (ii), (iii), (iv) is established at theorem level: (i) is an unproved phase-assignment convention, (ii) is an extension of a set-level projection to a phase-level statement not proved in~\citep{CascadeAlgSubstrate}, (iii) is a heuristic absorption, and (iv) presupposes the Kostant-gauge reduction of Hypothesis~\ref{hyp:kostant}. We record $\pi/87$ as the empirically-observed residual and flag its derivation as Open Item~5 in \S\ref{sec:status}.
\end{remark}

\begin{remark}[Recorded correspondence: $\alpha^{-1} \approx 137 + \pi/87$ under the stated identifications]\label{thm:alpha-full}
Combining the integer decomposition $137 = 87 + 50$ of Proposition~\ref{thm:alpha} with the unproved heuristic identification of the $\pi/87$ residue in Remark~\ref{rem:irrational}, the paper records the correspondence
\[
  \alpha^{-1} \;\approx\; 137 + \frac{\pi}{87} \;\approx\; 137.036\,110\,260.
\]
This is not a proved proposition about the physical constant. It is the combination of one proved arithmetic identity and one admitted heuristic identification, collected here as the $\alpha$-structural correspondence of the present paper. Five open items (B-formalisation, A-canonicality, C1 minimality, Kostant-gauge derivation, irrational-term derivation) are documented in \S\ref{sec:status}.
\end{remark}

\begin{remark}[Hypothesis audit]
The proof rests on: (i) Axiom $F1$; (ii) the C1 minimality principle; (iii) A's canonicality of the $\Zphi^2$-to-Coxeter-plane alignment; (iv) B's assignment of the upper exponents to the phason half; (v) P2's import of the $2I$ character table from Du~Val~\citep{DuVal1964}. Only (ii)--(iv) are not classical; (v) is an explicit external-literature import already catalogued in P2's scope section.
\end{remark}

\section{Numerical comparison}\label{sec:numerical}

\begin{table}[h]
\centering
\begin{tabular}{lll}
\toprule
Source & Value & Uncertainty \\
\midrule
CODATA~2018~\citep{CODATA2018} & $\alpha^{-1} = 137.035\,999\,084$ & $21 \times 10^{-9}$ \\
Remark~\ref{thm:alpha-full} (correspondence) & $\alpha^{-1} \approx 137 + \pi/87 \approx 137.036\,110\,260$ & --- \\
\midrule
Difference (ppm) & \multicolumn{2}{l}{$0.81$} \\
\bottomrule
\end{tabular}
\caption{Numerical agreement between the Remark~\ref{thm:alpha-full} structural correspondence and the CODATA~2018 recommended value for $\alpha^{-1}$. The correspondence is conditional on the full hypothesis load in \S\ref{sec:status}; the $0.81$-ppm agreement is a numerical consequence of the assumed identifications collected in this paper, not of a derived chain.}
\label{tab:numerical}
\end{table}

\section{Hypothesis status audit}\label{sec:status}

Table~\ref{tab:status} summarises the epistemic status of each ingredient.

\begin{table}[h]
\centering
\small
\begin{tabular}{p{3.3cm}p{6.5cm}p{5cm}}
\toprule
Layer & Content & Status \\
\midrule
$F1$ (\S\ref{sec:F1}) & $r = 1 + 1/r$, forcing $r = \varphi$ and $\Zphi$. & Axiom; classical once assumed. \\
Elser--Sloane (\S\ref{sec:F1}) & $E_8 = \{(x,\sigma(x)) : x \in \mathcal{I}\} \subset \mathcal{I} \oplus \sigma(\mathcal{I})$ via the icosian ring. & Theorem, classical. \\
C2 (\S\ref{sec:L12}) & Upward termination at $L_{12}$. & Proof attempt (Phase~B); status downgraded from ``classical'' to ``bespoke dimensional-counting argument documented in~\citep{CascadePhaseB}''. \\
C1 (\S\ref{sec:L12}) & Phason complement $\cong \Zphi^2$. & Conditional: closes under the $H_{\min}$ minimality principle \emph{plus} the admissibility hypotheses catalogued in~\citep{Cascade12D}. \\
C3 (\S\ref{sec:L12}) & $\Zphi$ self-reference. & Conditional on C1 plus classical $\Zphi$ structure; downstream, not independent. \\
A (\S\ref{sec:AB}) & $\Zphi^2$ aligns canonically with Coxeter planes. & Structurally argued; canonicality polish open. \\
B (\S\ref{sec:AB}) & Upper exponents $\{17,19,23,29\}$ identified with the phason sector, under B. & Argued from Galois halving; formal derivation open. \\
P2 \S 6 (\S\ref{sec:P2}) & $\C^{V_{600}} = W_{\mathrm{int}} \oplus W_{\mathrm{irr}} = 94 + 26$; $\{9,12,14,15\}$. & Imported computation (P2): computer-assisted shell/class substrate plus imported $2I$ character table from Du~Val. \\
Kostant gauge (\S\ref{sec:AB}) & $88 \to 87$ by $U(1) \times \Z/2$ gauge. & Open: a single global $U(1) \times \Z/2$ orbit removing exactly one slot is at present a convention, not a derivation. \\
Irrational-term (Remark~\ref{rem:irrational}) & $\pi / 87$ residue. & Open: structural identification, not a closed derivation. \\
\bottomrule
\end{tabular}
\caption{Hypothesis-status audit of Remark~\ref{thm:alpha-full}. No row claims theorem-grade closure from $F1$ alone. Five rows are flagged explicitly open.}
\label{tab:status}
\end{table}

\paragraph{What remains open for first-principles rigour.} Five items, in order of tightness of the remaining polish:
\begin{enumerate}
  \item \textbf{C1 minimality.} Reduce the ``$M$ is not a module over a ring strictly larger than $\Zphi$'' step of Phase~B to a classical theorem of algebraic number theory (candidate: uniqueness of the minimal ring containing $\varphi$ subject to the golden-ratio acceptance window).
  \item \textbf{A canonicality.} Formalise the ``canonical alignment of $\Zphi^2$ with the four Coxeter planes'' claim. Candidate: classification of $\sigma$-equivariant rank-$4$ sublattices inside $\R^8$.
  \item \textbf{B formalisation.} Lift the upper-exponents--phason identification from a Galois-halving argument to an intrinsic property of the $w$-eigenbasis.
  \item \textbf{Kostant-gauge derivation.} Justify $88 \to 87$ by deriving the one-slot reduction from an explicit $U(1) \times \Z/2$ gauge orbit on the phason data. The current paper treats this as a single shared-origin convention.
  \item \textbf{Irrational-term derivation.} Close the $\pi/87$ remainder by a derived chain (candidate: a phason-phase assignment $\pi$ per upper exponent, followed by a controlled two-to-one halving under $E_8 \to H_4$, followed by a derived $\Z/2$ absorption). All three steps are at present structural identifications; the modular arithmetic of a naive ``$88 \pi \to 44 \pi \to \text{residue} \to \pi/87$'' pipeline is not a proof.
\end{enumerate}
These five gaps do not change the arithmetic once assumed, but they prevent Identity~\eqref{eq:main} from being a theorem from $F1$ alone.

\section{Conclusion}

We have recorded a conditional structural correspondence $\alpha^{-1} \approx 137 + \pi/87$ under the hypothesis load $\{F1, \text{C1 (modulo minimality)}, A, B, \text{Kostant gauge}, P2\}$, plus the structural identification of the $\pi/87$ residue. (C2 and C3 belong to the ambient 12D closure narrative but do not enter the final arithmetic step.) Agreement with CODATA~2018 is $0.81$~ppm; the correspondence is a numerical consequence of the assumed identifications collected in this paper, not a derived chain. \textbf{Five open items} separate the present paper from a first-principles theorem: (i) C1 minimality; (ii) A canonicality; (iii) B formalisation; (iv) the Kostant-gauge reduction $88 \to 87$; and (v) the closed derivation of the $\pi/87$ residue.

\bibliographystyle{plainnat}
\begin{thebibliography}{99}

\bibitem[Cascade Foundations, 2026]{CascadeFoundations}
L.~Smart.
\newblock Cascade Foundations: the eight foundations F1--F8 underlying the {VFD} programme.
\newblock Preprint, Institute of Vibrational Field Dynamics, 2026.
\newblock Local source: \texttt{papers/paper-xxxvi/}.

\bibitem[Cascade {12D} Closure, 2026]{Cascade12D}
L.~Smart.
\newblock The {12D} closure theorem: phason complement uniqueness and self-reference on the upward cascade.
\newblock Preprint, Institute of Vibrational Field Dynamics, 2026.
\newblock Local source: \texttt{papers/cascade-12d-closure/cascade-12d-closure.tex}. Uses the $\Zphi^2$ convention for the phason complement, consistent with $\mathrm{rank}_\Z M = 4$. An earlier exploratory markdown note \texttt{papers/cascade-derivation/cascade-12d-closure.md} uses the older $\Zphi^4$ shorthand which is inconsistent with the stated Z-rank budget; the present paper follows the $\Zphi^2$ convention of the LaTeX preprint.

\bibitem[Cascade Phase {B}, 2026]{CascadePhaseB}
L.~Smart.
\newblock Phase {B}: closure of {C1, C2, C3}.
\newblock Preprint, Institute of Vibrational Field Dynamics, 2026.
\newblock Local source: \texttt{papers/cascade-derivation/cascade-phase-b-c1-c2-c3.md}.

\bibitem[Cascade {A}--{B} Closure, 2026]{CascadeABclosure}
L.~Smart.
\newblock Cascade assumptions {A} and {B}: closure of the canonical alignment and upper-exponent identifications.
\newblock Preprint, Institute of Vibrational Field Dynamics, 2026.
\newblock Local source: \texttt{papers/cascade-derivation/cascade-assumptions-a-b-closure.md}.

\bibitem[Cascade Identity-{C}, 2026]{CascadeIdentityC}
L.~Smart.
\newblock Cascade identity-{C} consolidated derivation.
\newblock Preprint, Institute of Vibrational Field Dynamics, 2026.
\newblock Local source: \texttt{papers/cascade-derivation/cascade-identity-c-consolidated.md}.

\bibitem[Cascade Algebraic Substrate, 2026]{CascadeAlgSubstrate}
L.~Smart.
\newblock Cascade--Classical Correspondence II: the cascade algebraic substrate ({P2}).
\newblock Preprint, Institute of Vibrational Field Dynamics, 2026.
\newblock Local source: \texttt{papers/cascade-algebraic-substrate/cascade-algebraic-substrate.tex}.

\bibitem[Smart, 2026 {V}]{SmartV}
L.~Smart.
\newblock Toward a spectral-geometric particle-mass model from the {600}-cell.
\newblock Preprint, Institute of Vibrational Field Dynamics, 2026.
\newblock Local source: \texttt{papers/paper-v/}.

\bibitem[Smart, 2026 {XXII}]{SmartXXII}
L.~Smart.
\newblock Paper {XXII}.
\newblock Preprint, Institute of Vibrational Field Dynamics, 2026.
\newblock Local source: \texttt{papers/paper-xxii/}.

\bibitem[Du~Val, 1964]{DuVal1964}
P.~Du~Val.
\newblock \emph{Homographies, Quaternions and Rotations}.
\newblock Oxford University Press, 1964.
\newblock Classical classification of finite subgroups of $\mathrm{Sp}(1)$; character table of $2I$ in Ch.~3.

\bibitem[Elser and Sloane, 1987]{ElserSloane1987}
V.~Elser and N.~J.~A. Sloane.
\newblock A highly symmetric four-dimensional quasicrystal.
\newblock \emph{Journal of Physics A: Mathematical and General}, 20(18):6161--6168, 1987.

\bibitem[Moody and Patera, 1993]{MoodyPatera1993}
R.~V. Moody and J.~Patera.
\newblock Quasicrystals and icosians.
\newblock \emph{Journal of Physics A: Mathematical and General}, 26(12):2829--2853, 1993.

\bibitem[Kostant, 1959]{Kostant1959}
B.~Kostant.
\newblock The principal three-dimensional subgroup and the {Betti} numbers of a complex simple {Lie} group.
\newblock \emph{American Journal of Mathematics}, 81:973--1032, 1959.

\bibitem[CODATA, 2021]{CODATA2018}
E.~Tiesinga, P.~J.~Mohr, D.~B.~Newell, and B.~N.~Taylor.
\newblock {CODATA} recommended values of the fundamental physical constants: 2018.
\newblock \emph{Reviews of Modern Physics}, 93:025010, 2021.

\end{thebibliography}

\end{document}
